\uuid{5dA2}
\titre{ Domaine, courbes de niveau et continuité }

\niveau{}
\module{ Fonctions de plusieurs variables }
\chapitre{Courbe de niveau}
\sousChapitre{}
\theme{}
\auteur{Quentin Liard}
\datecreate{ 2026-03-18}
\organisation{AMSCC}
\difficulte{3}

\contenu{

\texte{ 
On considère la fonction $f$ définie sur $\mathbb{R}^2$ par
$$
f(x,y)=\sqrt{4-x^2-y^2}.
$$
}

\begin{enumerate}

\item   \question{Déterminer le domaine de définition $D_f$ de $f$ et le représenter dans le plan. Cet ensemble est-il ouvert?}
\indication{}
\reponse{Pour que $f(x,y)$ soit définie, il faut
\[
4-x^2-y^2\geq 0,
\]
soit
\[
x^2+y^2\leq 4.
\]
Ainsi
\[
D_f=\{(x,y)\in\mathbb{R}^2\mid x^2+y^2\leq 4\}.
\]
C’est le disque fermé de centre $(0,0)$ et de rayon $2$.}

\item   \question{Déterminer les courbes de niveau de $f$, c’est-à-dire les ensembles $\{(x,y)\in D_f\mid f(x,y)=k\}$, selon les valeurs de $k$.}
\indication{}
\reponse{On cherche les points tels que
\[
\sqrt{4-x^2-y^2}=k.
\]
Cela impose d’abord $k\geq 0$. En élevant au carré :
\[
4-x^2-y^2=k^2
\quad\Longleftrightarrow\quad
x^2+y^2=4-k^2.
\]
Donc :
\begin{itemize}
\item si $0\leq k\leq 2$, la courbe de niveau est le cercle de centre $(0,0)$ et de rayon $\sqrt{4-k^2}$ ;
\item si $k<0$ ou $k>2$, il n’y a pas de courbe de niveau.
\end{itemize}}

\item   \question{La fonction $f$ est-elle continue sur son domaine de définition $D_f$ ?}
\indication{}
\reponse{Oui. La fonction $(x,y)\mapsto 4-x^2-y^2$ est polynomiale donc continue sur $\mathbb{R}^2$, et la fonction racine carrée est continue sur $[0,+\infty[$. Par composition, $f$ est continue sur son domaine $D_f$.}

\item   \question{Calculer les dérivées partielles $\dfrac{\partial f}{\partial x}$ et $\dfrac{\partial f}{\partial y}$ sur l'intérieur de $D_f$.}
\indication{}
\reponse{Pour $x^2+y^2<4$, on a
\[
f(x,y)=(4-x^2-y^2)^{1/2}.
\]
Donc
\[
\frac{\partial f}{\partial x}(x,y)=\frac12(4-x^2-y^2)^{-1/2}(-2x)
=\frac{-x}{\sqrt{4-x^2-y^2}},
\]
et de même
\[
\frac{\partial f}{\partial y}(x,y)=\frac{-y}{\sqrt{4-x^2-y^2}}.
\]}
\item \question{Les dérivées partielles en $(2,0)$ et $(0,2)$ existent-elles?}

\reponse{
On étudie l'existence des dérivées partielles de
\[
f(x,y)=\sqrt{4-x^2-y^2}
\]
aux points \((2,0)\) et \((0,2)\).

\medskip

\textbf{En \((2,0)\).} On a \(f(2,0)=0\).

Pour la dérivée partielle par rapport à \(x\), on considère
\[
\frac{\partial f}{\partial x}(2,0)
=\lim_{h\to 0}\frac{f(2+h,0)-f(2,0)}{h},
\]
si cette limite existe. Or, pour que \((2+h,0)\in D_f\), il faut
\[
(2+h)^2\le 4,
\]
donc \(h\le 0\) au voisinage de \(0\). Pour \(h<0\),
\[
f(2+h,0)=\sqrt{4-(2+h)^2}=\sqrt{-4h-h^2}.
\]
Ainsi,
\[
\frac{f(2+h,0)-f(2,0)}{h}
=\frac{\sqrt{-4h-h^2}}{h}.
\]
Quand \(h\to 0^{-}\), ce quotient tend vers \(-\infty\). La dérivée partielle
\[
\frac{\partial f}{\partial x}(2,0)
\]
n'existe donc pas.

Pour la dérivée partielle par rapport à \(y\), on regarde
\[
\frac{\partial f}{\partial y}(2,0)
=\lim_{h\to 0}\frac{f(2,\,h)-f(2,0)}{h}.
\]
Mais
\[
f(2,h)=\sqrt{4-4-h^2}=\sqrt{-h^2},
\]
ce qui n'est défini que pour \(h=0\). Il n'existe donc pas de voisinage de \(0\) dans lequel ce quotient soit défini pour \(h\neq 0\). Par conséquent,
\[
\frac{\partial f}{\partial y}(2,0)
\]
n'existe pas.

\medskip

\textbf{En \((0,2)\).} On a \(f(0,2)=0\).

Pour la dérivée partielle par rapport à \(x\), on considère
\[
\frac{\partial f}{\partial x}(0,2)
=\lim_{h\to 0}\frac{f(h,2)-f(0,2)}{h}.
\]
Or
\[
f(h,2)=\sqrt{4-h^2-4}=\sqrt{-h^2},
\]
qui n'est défini que pour \(h=0\). Donc
\[
\frac{\partial f}{\partial x}(0,2)
\]
n'existe pas.

Pour la dérivée partielle par rapport à \(y\), on considère
\[
\frac{\partial f}{\partial y}(0,2)
=\lim_{h\to 0}\frac{f(0,2+h)-f(0,2)}{h}.
\]
Pour que \((0,2+h)\in D_f\), il faut
\[
(2+h)^2\le 4,
\]
donc \(h\le 0\) au voisinage de \(0\). Pour \(h<0\),
\[
f(0,2+h)=\sqrt{4-(2+h)^2}=\sqrt{-4h-h^2}.
\]
Ainsi,
\[
\frac{f(0,2+h)-f(0,2)}{h}
=\frac{\sqrt{-4h-h^2}}{h},
\]
et ce quotient tend vers \(-\infty\) quand \(h\to 0^{-}\). Donc
\[
\frac{\partial f}{\partial y}(0,2)
\]
n'existe pas.

\medskip

En conclusion, en \((2,0)\) et en \((0,2)\), les dérivées partielles de \(f\) n'existent pas.
}
\end{enumerate}

}